%% =============================================================================
%%  Etapa 6 — Diálogo com a tradição
%% =============================================================================
\chapter{Diálogo}\label{cap:dialogo}

\begin{objetivo}
Pôr as observações das Etapas 1–5 em diálogo com outros leitores: comentários
exegéticos, obras de referência e a história da interpretação. O objetivo não
é encontrar \enquote{a resposta}, mas testar as próprias conclusões e
descobrir o que passou despercebido.
\end{objetivo}

\section{Como consultar comentários}

Leia cada comentário com as perguntas das etapas anteriores na mão. Anote
onde ele confirma, corrige ou amplia o que você viu e, principalmente,
\emph{que argumentos} usa. Comentários de tradições diferentes mostram
pressupostos diferentes; ler mais de um é a melhor proteção contra ler só o
que já se pensava \parencite{fee2002, osborne2006}.

\begin{ebtabela}{@{}>{\raggedright}p{.26\linewidth} >{\raggedright}p{.34\linewidth} >{\raggedright\arraybackslash}X@{}}
\EBcab{Tipo} & \EBcab{Exemplos para Lucas} & \EBcab{Quando usar}\\\midrule
Técnico (texto grego) & \textcite{bovon2002}; \textcite{fitzmyer1981}; \textcite{marshall1978}; \textcite{nolland1989} & crítica textual, gramática, história da tradição\\
Exegético (tradução) & \textcite{bock1994}; \textcite{green1997} & argumento do texto, contexto narrativo e social\\
Uso do AT & \textcite{beale2007} (Pao e Schnabel, \enquote{Luke}) & citações, alusões, ecos\\
Léxicos e gramáticas & \textcite{louwnida1988}; \textcite{runge2010} & campos semânticos, marcadores de discurso\\
Crítica textual & \textcite{metzger1994} & variantes discutidas pelo comitê da UBS\\
\end{ebtabela}

\begin{nota}[Em português]
\textcite{wegner1998} e \textcite{cassio2000} são manuais de exegese com
exercícios; \textcite{egger1994} apresenta os métodos sincrônicos (análise
literária e narrativa); \textcite{feestuart2014} é uma boa introdução geral.
\end{nota}

\section{Ficha de leitura}

Para cada obra consultada, preencha:

\begin{ebtabela}{@{}>{\raggedright}p{.3\linewidth} >{\raggedright\arraybackslash}X@{}}
\EBcab{Campo} & \EBcab{O que registrar}\\\midrule
Referência & autor, obra, páginas\\
Tese sobre o trecho & em uma frase\\
Argumentos & textuais, gramaticais, históricos, teológicos\\
Concorda comigo em & \ldots\\
Discorda de mim em & \ldots\ e por quê\\
Pergunta nova & algo que eu não tinha visto\\
\end{ebtabela}

\section{Questões em debate}

Algumas perguntas que os comentários discutem e que este estudo já tocou:

\paragraph{Quem são os pobres?} Pobres materiais, os \enquote{pobres de
YHWH} dos Salmos (piedosos e dependentes de Deus), ou os discípulos que
deixaram tudo (\rb{Lc 5:11, 28})? Compare as respostas de pelo menos dois
comentários com a sua resposta à pergunta~\ref{perg:pobres}.

\paragraph{Promessa ou ética?} As felicidades descrevem uma virtude a
buscar ou anunciam uma ação de Deus? A análise da Etapa~3 (macarismo
escatológico) e do passivo divino (Etapa~2) pesa a favor de qual leitura?

\paragraph{Os ais: condenação ou advertência?} O \enquote{ai} fecha a porta
aos ricos ou os chama à mudança, como Zaqueu (\rb{19:1–10})?

\section{História da recepção}

Ler como o texto foi lido ao longo dos séculos mostra possibilidades e
limites de interpretação. Um exemplo: Ambrósio de Milão, no século~IV,
associa as quatro felicidades de Lucas às quatro virtudes cardeais
(temperança, justiça, prudência e fortaleza) em sua \emph{Exposição do
Evangelho segundo Lucas} (livro~V). A leitura é engenhosa, mas mostra
como uma estrutura filosófica pode ser projetada sobre o texto.

\begin{pergunta}
Que pressupostos da sua própria tradição de leitura influenciam o modo como
você lê \enquote{pobres} e \enquote{ricos}? Onde o texto confirma esses
pressupostos e onde os questiona?
\end{pergunta}
